\section{Verifier y LLM-as-a-Judge: incorporar la evaluación al razonamiento}

\subsection{De la evaluación al final a la evaluación intermedia}
Si el éxito solo se determina al final de la trayectoria, la retroalimentación llega demasiado tarde. El ponente propone adelantar el verifier a los pasos intermedios, de modo que el sistema pueda juzgar, durante el razonamiento mismo, si «vale la pena continuar por el camino actual».

\begin{importantbox}{Las tres funciones del verifier}
\begin{enumerate}
    \item Puntuación de rutas: sirve de base para podar la búsqueda;
    \item Detección de errores: identifica estados que se han desviado de forma evidente;
    \item Filtrado de datos: selecciona trayectorias de alta calidad utilizables para el entrenamiento.
\end{enumerate}
\end{importantbox}

\subsection{El valor práctico de LLM-as-a-Judge}
Hacer que el LLM evalúe la trayectoria suele ser más estable que hacer que genere directamente la trayectoria óptima. Esto se debe a que «juzgar si algo es correcto» suele ser más fácil que «acertar de un solo paso», una lección práctica que este apartado reitera en varias ocasiones.

\begin{knowledgebox}{Elementos clave del prompt del juez}
\begin{itemize}
    \item El objetivo y las restricciones de la tarea;
    \item La trayectoria completa (o un resumen de los pasos clave);
    \item Una plantilla de salida explícita: aprobado o rechazado, junto con la justificación;
    \item Una estimación de confianza, que la estrategia de búsqueda pueda usar.
\end{itemize}
\end{knowledgebox}

\subsection{Confiabilidad y calibración del juez}
El juez no constituye una verdad absoluta. Si el criterio de evaluación es ambiguo o el prompt se ha desviado de su propósito original, el juez puede llegar a una conclusión errónea con alta confianza y, con ello, desviar la búsqueda.

\begin{warningbox}{Tres riesgos del juez}
\begin{itemize}
    \item Manipulación de la recompensa: el modelo aprende a complacer el estilo del juez en lugar de completar la tarea de verdad;
    \item Sesgo de dominio: el juez pierde precisión ante plantillas de páginas web con las que no está familiarizado;
    \item Desajuste de confianza: una puntuación alta no implica una tasa de éxito alta.
\end{itemize}
\end{warningbox}

\begin{figure}[H]
\centering
\includegraphics[width=0.82\textwidth]{frame-08-full-pipeline.jpg}
\caption{Minuto aproximado 01:12:10 del video: la planificación, la ejecución y el verificador forman un ciclo cerrado completo; la evaluación ya no ocurre únicamente al final de la tarea}
\end{figure}

\subsection{Resumen de la sección}
El verifier hace que el agente pase de «actuar primero y juzgar después» a «juzgar mientras actúa». Sin embargo, la confiabilidad del juez requiere una calibración continua, pues de lo contrario puede orientar al sistema hacia una dirección de optimización equivocada.
